\chapter{Adenoidectomy}

\section*{What Is This Test or Procedure?}
Adenoidectomy is an operation to remove enlarged or chronically infected adenoids, lymphoid tissue located behind the nose. It is performed most often in children.

\section*{Alternative Names}
Adenoid removal, adenoid surgery

\section*{Why This Test or Procedure Is Ordered}
It may be recommended for persistent nasal obstruction, sleep-disordered breathing or obstructive sleep apnea, recurrent ear disease, or chronic adenoid infection when medical treatment is insufficient.

\section*{How This Test or Procedure Is Performed}
Under general anesthesia, the surgeon removes the adenoids through the mouth using a curette, powered instrument, cautery, or another technique. Tonsillectomy or ear-tube placement may be performed during the same anesthetic when indicated.

\section*{Preparation, Risks, and Recovery}
Follow fasting and medication instructions and report bleeding problems or loose teeth. Risks include bleeding, infection, dehydration, nausea, temporary sore throat, voice or swallowing changes, dental injury, and regrowth. Most patients go home the same day; fluids, pain control, and follow-up are important.

\section*{References}
\begin{enumerate}
  \item NHS. Adenoidectomy. 2025.
  \item MedlinePlus. Enlarged Adenoids and Adenoidectomy. U.S. National Library of Medicine; 2025.
\end{enumerate}

\clearpage
